% -------------------------------------------------------------------------
% WBS ID: RES-GEO-MET
% Deliverable: Geometric and Hydraulic Descriptor Framework
% Status: In progress
% Evidence status: Minimum descriptors tied to hydrodynamic outputs
% Review status: Not reviewed
% Last updated: 2026-07-19
% Dependencies: RES-GEO-PAR, RES-DAT-SCN
% Downstream interfaces: RES-PHY-BAR, RES-PHY-LOD, RES-VER-MET, Software output schema
% -------------------------------------------------------------------------

\section{RES-GEO-MET --- Geometric and Hydraulic Descriptor Framework}
\label{sec:res-geo-met}

\subsection{Purpose and Scope}
\label{sec:res-geo-met-purpose}

This Deliverable defines the minimum geometric and hydraulic
descriptors needed to compare wall-type and obstacle/array-type
barriers in fixed-rigid local tsunami simulations. Descriptor records
shall preserve provenance and uncertainty. Damage state, material
deformation and scour descriptors are recorded only as deferred
structural extensions.

\subsection{Research Questions}
\label{sec:res-geo-met-research-questions}

\begin{enumerate}
  \item Which descriptors are required before local hydrodynamic
    comparison can begin?
  \item Which descriptors drive software geometry, and which are
    post-processing metrics emitted by the solver?
  \item Which descriptors must carry uncertainty/provenance flags?
\end{enumerate}

\subsection{Terminology and Definitions}
\label{sec:res-geo-met-definitions}

% TODO: Define the technical terminology, symbols, quantities and
% classifications used in this Deliverable.

\subsection{Literature Evidence}
\label{sec:res-geo-met-literature-evidence}

\textcite{deljesus2012porous} links vertical barrier form, porosity
and transmission to hydraulic response. Source role: supports
including blockage, opening and transmission descriptors.
\textcite{watanabe2022validation} links local geometry to separate
force, pressure, velocity and free-surface validation quantities.
Project conclusion: descriptors must identify both what was simulated
and which hydrodynamic outputs are meaningful for comparison.

\subsection{Governing Relations and Quantitative Definitions}
\label{sec:res-geo-met-governing-relations}

For array and obstacle cases, the plan blockage ratio shall be
reported as:

\begin{align}
  \beta_{\mathrm{blk}}
  &=
  \frac{
    A_{\perp,\mathrm{solid}}
  }{
    A_{\perp,\mathrm{corridor}}
    +
    \epsilon_A
  }
  \label{eq:res-geo-met-blockage-ratio}
\end{align}

where $A_{\perp,\mathrm{solid}}$ is the projected solid area normal to
the corridor centreline and $A_{\perp,\mathrm{corridor}}$ is the
available projected corridor area at the comparison section. Decision
role: \cref{eq:res-geo-met-blockage-ratio} is a comparison descriptor
for wall/obstacle/array cases. Supporting evidence:
\textcite{deljesus2012porous}. Limitation: it does not replace full
3D flow solution because wake, overtopping and pressure depend on
local geometry and bathymetry.

\subsection{Geometry Family or Representation}
\label{sec:res-geo-met-geometry-family-representation}

% TODO: Complete this RES-GEO Domain-specific subsection.

\subsection{Design Variables and Parameter Bounds}
\label{sec:res-geo-met-design-variables-parameter-bounds}

% TODO: Complete this RES-GEO Domain-specific subsection.

\subsection{Geometric Descriptors}
\label{sec:res-geo-met-geometric-descriptors}

Required pre-simulation descriptors are class label, local corridor
coordinates, orientation, crest elevation, height, thickness or
element width, footprint length, spacing, row count, opening fraction,
roughness flag, blockage ratio and provenance quality.

Required post-simulation hydraulic descriptors are maximum run-up,
overtopping volume or discharge, transmitted and reflected energy
measures, downstream maximum depth and velocity, peak pressure,
resultant force, impulse and moment. Damage, survival state and
structural-capacity proxies are retained only for later RES-STR
screening and shall not be used to claim baseline hydrodynamic
performance.

\subsection{Placement and Arrangement Rules}
\label{sec:res-geo-met-placement-arrangement-rules}

% TODO: Complete this RES-GEO Domain-specific subsection.

\subsection{Feasibility Constraints}
\label{sec:res-geo-met-feasibility-constraints}

% TODO: Complete this RES-GEO Domain-specific subsection.

\subsection{Two- and Three-Dimensional Representation}
\label{sec:res-geo-met-two-three-dimensional-representation}

% TODO: Complete this RES-GEO Domain-specific subsection.

\subsection{Candidate Approaches}
\label{sec:res-geo-met-candidate-approaches}

% TODO: Define the credible candidate approaches considered.

\subsection{Critical Comparison}
\label{sec:res-geo-met-critical-comparison}

% TODO: Compare candidates using physically and project-relevant criteria.
% Do not select an approach solely on popularity or implementation ease.

\subsection{Assumptions and Applicability}
\label{sec:res-geo-met-assumptions}

% TODO: State assumptions and the conditions under which they are valid.

\subsection{Limitations and Failure Conditions}
\label{sec:res-geo-met-limitations}

% TODO: State where the evidence, model or selected approach ceases to be
% reliable or applicable.

\subsection{Project-Specific Conclusions}
\label{sec:res-geo-met-conclusions}

The geometry descriptor framework is sufficient when each barrier case
can be reconstructed, meshed, run and post-processed using stable
identifiers. It is not sufficient for structural certification or
damage prediction.

\subsection{Selected Approach and Rationale}
\label{sec:res-geo-met-selected-approach}

The selected descriptor approach separates pre-simulation geometry
fields from solver-emitted hydrodynamic metrics. This keeps software
data handling parallel with Research without moving solver
implementation into RES-GEO.

\subsection{Unresolved Decisions}
\label{sec:res-geo-met-unresolved-decisions}

Open decisions are the final blockage-normalisation section, roughness
descriptor values, gap/opening representation and accepted uncertainty
classes for inferred geometry.

\subsection{Requirements and Handoffs}
\label{sec:res-geo-met-requirements}

Software shall emit descriptors and metrics using the same
\texttt{corridor\_id}, \texttt{geometry\_id} and \texttt{case\_id}.
RES-PHY-BAR and RES-PHY-LOD shall define the hydrodynamic meaning of
the emitted metrics; RES-VER-MET shall define acceptance tolerances.

\subsection{Deliverable Completion Record}
\label{sec:res-geo-met-completion-record}

% TODO: Record whether the Deliverable is:
%
% - Not started
% - In progress
% - Draft complete
% - Reviewed
% - Accepted
%
% Acceptance requires:
%
% - the research questions to be answered;
% - claims to be supported by appropriate evidence;
% - assumptions and limitations to be explicit;
% - the project-specific conclusion to be stated;
% - downstream requirements to be traceable;
% - unresolved decisions to be closed or formally deferred.
